\chapter{因果排序与邻接矩阵}
\label{app:orderings}
\begin{definition}
给定有向无环图 $\mathcal G$，如果一个排列，即双射
\[\pi\colon\{1,\ldots,p\}\to\{1,\ldots,p\},\]
满足
\[\pi(i)<\pi(j)\qquad\text{只要 }j\in\DE_i^{\mathcal G},\]
则称其为因果排序（causal ordering），有时也称拓扑排序（topological ordering）。
\end{definition}
由于 DAG 具有无环结构，总存在拓扑排序，见命题 B.2。但这一排序未必唯一。节点 $\pi^{-1}(1)$ 没有父节点，因此是源节点（source node）；节点 $\pi^{-1}(p)$ 没有后代节点，因此是汇节点（sink node）。

\begin{proposition}每个 DAG 都存在拓扑排序。\end{proposition}
\begin{proof}
采用归纳法。只需证明每个 DAG 中都存在没有祖先的节点。从任意节点出发，若它有父节点，就移动到其中一个父节点。这个过程中不可能再次访问已经访问过的节点，否则就会构成有向环。至多经过 $p-1$ 步，就会到达一个没有父节点的节点。
\end{proof}

\begin{definition}
对于有 $d$ 个节点的有向图 $\mathcal G=(\mathcal V,\mathcal E)$，可以用一个元素取值为 $0$ 或 $1$ 的 $d\times d$ 二元矩阵 $A$ 表示：
\[A_{i,j}=1\quad\Longleftrightarrow\quad(i,j)\in\mathcal E.\]
称 $A$ 为 $\mathcal G$ 的邻接矩阵（adjacency matrix）。
\end{definition}

DAG 的这种表示特别适合高效实现算法。关于邻接矩阵的运算，有若干有用结果，下面列举其中一部分。
\begin{remark}
\begin{enumerate}[label=(\roman*)]
\item 设 $A$ 是 DAG $\mathcal G$ 的邻接矩阵。矩阵平方 $A^2$ 的第 $(i,j)$ 个元素，等于从 $i$ 到 $j$ 的长度为二的路径数。这是因为
\[(A^2)_{i,j}=\sum_k A_{ik}A_{kj}.\]
\item 一般地，
\[(A^k)_{ij}=\#\{\text{从 }i\text{ 到 }j\text{、长度为 }k\text{ 的路径}\}.\]
\item 如果沿有向路径节点索引递增，即 $j\in\DE_i^{\mathcal G}$ 蕴含 $j>i$，那么恒等映射就是一个因果排序，且邻接矩阵为上三角矩阵，即非零元素只可能出现在右上半部。
\item 若图是稀疏的，可以采用稀疏矩阵，以节省存储空间和计算时间。
\end{enumerate}
\end{remark}

Robinson [1970, 1973] 以及独立开展研究的 Stanley [1973]，都研究了包含 $d$ 个节点的 DAG 的数量。这类矩阵或 DAG 的数量随 $d$ 增长得非常快，见表 B.1。

McKay [2004] 证明了 DAG 的如下等价刻画；这一结果此前由 Eric W. Weisstein 作为猜想提出。
\begin{theorem}
矩阵 $A$ 是某个 DAG $\mathcal G$ 的邻接矩阵，当且仅当 $A+I_d$ 是一个所有特征值均为实数且严格大于零的 $0$--$1$ 矩阵。
\end{theorem}

\begin{table}[!htbp]
\centering
\small
\begin{tabular}{r l}
\toprule
$d$ & 含 $d$ 个节点的 DAG 数量 \\
\midrule
1 & 1 \\
2 & 3 \\
3 & 25 \\
4 & 543 \\
5 & 29281 \\
6 & 3781503 \\
7 & 1138779265 \\
8 & 783702329343 \\
9 & 1213442454842881 \\
10 & 4175098976430598143 \\
11 & 31603459396418917607425 \\
12 & 521939651343829405020504063 \\
13 & 18676600744432035186664816926721 \\
14 & 1439428141044398334941790719839535103 \\
15 & 237725265553410354992180218286376719253505 \\
16 & 83756670773733320287699303047996412235223138303 \\
17 & 62707921196923889899446452602494921906963551482675201 \\
18 & 99421195322159515895228914592354524516555026878588305014783 \\
19 & 332771901227107591736177573311261125883583076258421902583546773505 \\
\bottomrule
\end{tabular}
\caption{DAG 的数量与节点数 $d$ 的关系，数据取自 \url{http://oeis.org/A003024} [OEIS Foundation Inc., 2017]。这些数的位数增长得比任意线性项都快。}
\end{table}
